\scanpage{16}
% Mirror-image typescript on the other side.
$G$ gpe de Lie compact connexe

\diagram{\begin{tikzcd}[ampersand replacement=\&,column sep=small,row sep=small]
C_{(0)}^{p,q}\arrow[d,phantom,"\cup" description] \& {}=\Omega^q(\unsure{\mathord{-}}G^{p+1}/G) \&\\
C_{(1)}^{pq}\arrow[d,phantom,"\cup" description] \& {}=\omega^{(q)}(G^{p+1}/G) \&\\
C_{(2)}^{pq} \& {}=\omega_{\mathrm{inv}}^{(q)}(G^{p+1}/G) \&
\end{tikzcd}}
\margin{inclusion induit un isom sur \unsure{les termes|le terme|les dernières} $\,{}'E_1$ (cohomologie pour $p$ fixe}
\margin{id}

Dans ce dernier bicomplexe, l'opérateur différentiel par rapport à $q$ \uline{est nul}. Donc on trouve
\[
\begin{aligned}
H^*(B_G,\mathbb R)&=H^*(C_{\bullet}^{pq})\simeq H^*(C_{(1)}^{pq})
\simeq\sum_{p,q}\bigl\{H^p\bigl(\unsure{\nu|v}\longrightarrow
\omega_{\mathrm{inv}}^{(q)}(G^{\nu+1}/\unsure{G})\unsure{)\}|)}\\[-2pt]
&\hspace{15em}\rotatebox{90}{$=$}\\
&\hspace{15em}E_2^{pq}
\end{aligned}
\]
(isomorphismes canoniques)

\uuline{La suite spectrale dégénère} : différentielles $d_r$ nulles pour $r\geq2$.

Soit $T\subset\omega^*_{\mathrm{inv}}(G)$, \unsure{élémts|éléments} transgressifs. (\unsure{Ils|OK} sont de degrés\unsure{$+$|s} impairs)

On a
\[
T\xrightarrow{\varphi} E_2^{1*}=\omega_{\mathrm{inv}}^*(G\times G/G)
\subset\underset{\substack{\rotatebox{90}{$=$}\\
  \omega_{\mathrm{inv}}^*(G)\otimes\omega_{\mathrm{inv}}^*(G)}}
{\omega_{\mathrm{inv}}^*(G\times G)}
\]
par
\[
\varphi(T)=T\otimes1-1\otimes T
\]
\uline{Th} . $\mathrm{Sym}^*(T)\longrightarrow H^*(B_G,\mathbb R)$ \uline{déduit} de $\varphi$ \uline{est} \uline{un isomorphisme}.
% archive sheet number: 15
